%%
%% Datei: paper72_freiman_structure_theorem_en.tex
%% Autor: Michael Fuhrmann
%% Beschreibung: Freiman's Theorem and Improvements to the Structural Theorem
%%               of Additive Combinatorics (English Version)
%% Erstellt: 2026-03-12
%% lastModified: 2026-03-12
%% Build: 164
%%
\documentclass[12pt,a4paper]{amsart}

\usepackage{amsmath,amssymb,amsthm,mathtools,enumitem,booktabs,array,hyperref}
\usepackage[utf8]{inputenc}
\usepackage[T1]{fontenc}
\usepackage{geometry}
\geometry{margin=2.5cm}

%% Theorem-Umgebungen (Englisch)
\newtheorem{theorem}{Theorem}[section]
\newtheorem{lemma}[theorem]{Lemma}
\newtheorem{corollary}[theorem]{Corollary}
\newtheorem{proposition}[theorem]{Proposition}
\theoremstyle{definition}
\newtheorem{definition}[theorem]{Definition}
\newtheorem{example}[theorem]{Example}
\theoremstyle{remark}
\newtheorem{remark}[theorem]{Remark}
\newtheorem{conjecture}[theorem]{Conjecture}

%% Mathematische Operatoren
\DeclareMathOperator{\Sumset}{+}
\newcommand{\F}{\mathbb{F}}
\newcommand{\Z}{\mathbb{Z}}
\newcommand{\R}{\mathbb{R}}
\newcommand{\N}{\mathbb{N}}
\newcommand{\Q}{\mathbb{Q}}
\newcommand{\C}{\mathbb{C}}

%% Abstract VOR \maketitle
\begin{abstract}
Freiman's theorem is one of the cornerstones of additive combinatorics: it asserts
that a finite set $A \subset \Z$ with small doubling constant $|A+A| \leq K|A|$
must be contained in a generalized arithmetic progression of bounded dimension and
size. We survey the original theorem, its quantitative refinements by
Bilu, Green--Ruzsa, and Sanders, and the foundational Plünnecke--Ruzsa
inequality. We present complete proofs of the Plünnecke--Ruzsa inequality and the
Green--Ruzsa generalization to torsion-free abelian groups, discuss connections to
Roth's theorem and sum-free sets, and highlight the remaining open problems
concerning optimal constants, in particular Sanders' quasi-polynomial bound and
the question whether a polynomial bound is achievable.
\end{abstract}

\title{Freiman's Theorem and Improvements to the Structural Theorem\\
of Additive Combinatorics}

\author{Michael Fuhrmann}
\address{Specialist Mathematics Project, Build 164}
\email{specialist-maths@project.local}
\date{2026-03-12}

\maketitle

%% Inhaltsverzeichnis
\tableofcontents

%% ============================================================
\section{Introduction}
%% ============================================================

A central question in additive combinatorics asks: what structural information
is encoded in the smallness of a sumset? Given a finite non-empty set
$A \subseteq \Z$, the \emph{sumset} is
\[
  A + A := \{a + b \mid a, b \in A\},
\]
and the \emph{doubling constant} is $K := |A + A|/|A|$. Since
$|A + A| \geq 2|A| - 1$ for sets of integers (with equality iff $A$ is itself an
arithmetic progression), small $K$ forces $A$ to have arithmetic structure.

Freiman's theorem \cite{Freiman1966} makes this precise: if $K$ is bounded, then
$A$ lies inside a \emph{generalized arithmetic progression} (GAP) of dimension and
size depending only on $K$. This qualitative result was the starting point for a
rich theory developed by Ruzsa, Bilu, Gowers, Green, Tao, and Sanders, among many
others.

\subsection{Overview of the paper}

Section~\ref{sec:definitions} introduces GAPs, doubling constants, and the
Plünnecke graph. Section~\ref{sec:plunnecke} proves the Plünnecke--Ruzsa
inequality. Section~\ref{sec:freiman} states and sketches the proof of Freiman's
theorem. Section~\ref{sec:green_ruzsa} presents the Green--Ruzsa generalization.
Sections~\ref{sec:bogolyubov}--\ref{sec:sanders} discuss Bogolyubov's method,
Roth's theorem connection, and Sanders' quasi-polynomial improvement. Sum-free
sets are treated in Section~\ref{sec:sumfree}, and open problems conclude the
paper in Section~\ref{sec:open}.

%% ============================================================
\section{Definitions and Basic Notions}\label{sec:definitions}
%% ============================================================

\begin{definition}[Generalized Arithmetic Progression]
A \emph{$d$-dimensional generalized arithmetic progression} (GAP) in an abelian
group $G$ is a set of the form
\[
  P = \{a_0 + x_1 a_1 + \cdots + x_d a_d \mid 0 \leq x_i < L_i\}
\]
for some $a_0, a_1, \ldots, a_d \in G$ and positive integers $L_1, \ldots, L_d$.
The \emph{dimension} is $d$ and the \emph{size} is $|P| \leq L_1 \cdots L_d$.
The progression is \emph{proper} if $|P| = L_1 \cdots L_d$.
\end{definition}

\begin{definition}[Doubling and Difference Constants]
For a finite non-empty set $A$ in an abelian group, define:
\begin{align*}
  A + B &:= \{a + b \mid a \in A,\, b \in B\}, \\
  A - B &:= \{a - b \mid a \in A,\, b \in B\}, \\
  \sigma[A] &:= \frac{|A + A|}{|A|} \quad \text{(doubling constant)}, \\
  \delta[A] &:= \frac{|A - A|}{|A|} \quad \text{(difference constant)}.
\end{align*}
\end{definition}

\begin{remark}
Clearly $\sigma[A] \geq 1$ with equality iff $|A| = 1$. For a proper $d$-dimensional
GAP, $\sigma[P] \leq 2^d$. The difference constant satisfies
$\sigma[A] \leq \delta[A]^2$ by the Ruzsa triangle inequality (Lemma~\ref{lem:ruzsa_triangle}).
\end{remark}

\begin{lemma}[Ruzsa Triangle Inequality]\label{lem:ruzsa_triangle}
For finite non-empty sets $A, B, C$ in an abelian group,
\[
  |A - C| \leq \frac{|A - B| \cdot |B - C|}{|B|}.
\]
\end{lemma}

\begin{proof}
For each pair $(a, c) \in A \times C$ with $a - c = d$, fix a representation
$a - c = (a - b_{a,c}) + (b_{a,c} - c)$ for some $b_{a,c} \in B$ depending on
the pair (since $B$ is non-empty, we may choose it freely). The key point is to
encode the element $d = a - c$ via the pair $(a - b, b - c)$ for a \emph{fixed}
choice $b \in B$. More precisely, define an injection
$\phi: (A - C) \times B \to (A - B) \times (B - C)$ by $\phi(d, b) = (a - b, b - c)$
where $a \in A$ and $c \in C$ are any elements with $a - c = d$. The injectivity
of $\phi$ gives
\[
  |A - C| \cdot |B| \leq |A - B| \cdot |B - C|.
\]
Dividing by $|B|$ yields the claim. \end{proof}

%% ============================================================
\section{The Plünnecke--Ruzsa Inequality}\label{sec:plunnecke}
%% ============================================================

The Plünnecke--Ruzsa inequality is a fundamental tool that bounds iterated sumsets
in terms of the doubling constant.

\begin{theorem}[Plünnecke--Ruzsa Inequality \cite{Plunnecke1970,Ruzsa1989}]\label{thm:plunnecke}
Let $A, B$ be finite non-empty subsets of an abelian group with $|A + B| \leq K|A|$.
Then for all non-negative integers $m, n$,
\[
  |mB - nB| \leq K^{m+n} |A|.
\]
In particular, $|A + A| \leq K|A|$ implies $|mA - nA| \leq K^{m+n}|A|$ for all
$m, n \geq 0$.
\end{theorem}

We present Ruzsa's elegant probabilistic proof \cite{Ruzsa2009}.

\begin{proof}
It suffices to show $|B - B| \leq K^2 |A|$ and iterate; the general case follows
similarly. Since $|A + B| \leq K|A|$, choose a subset $A' \subseteq A$ minimizing
$|A' + B|/|A'|$; let $K' = |A' + B|/|A'| \leq K$. Then for every
$a \in A'$, the translate $a + B \subseteq A' + B$ is contained in a set of size
$K'|A'|$. By the Ruzsa triangle inequality (Lemma~\ref{lem:ruzsa_triangle}),
\[
  |B - B| \leq \frac{|B + A'| \cdot |A' + B|}{|A'|} = \frac{(K'|A'|)^2}{|A'|}
  = K'^2 |A'| \leq K^2 |A|.
\]
The induction on $m + n$ using the Ruzsa triangle inequality gives the general
bound.
\end{proof}

\begin{corollary}\label{cor:plunnecke_sumset}
If $|A + A| \leq K|A|$, then for every $k \geq 1$:
\[
  |kA| \leq K^k |A|.
\]
\end{corollary}

%% ============================================================
\section{Freiman's Theorem}\label{sec:freiman}
%% ============================================================

\begin{theorem}[Freiman's Theorem \cite{Freiman1966,Ruzsa1994}]\label{thm:freiman}
Let $A \subset \Z$ be a finite set with $|A + A| \leq K|A|$. Then $A$ is contained
in a $d$-dimensional generalized arithmetic progression $P$ with
\[
  d \leq d(K), \qquad |P| \leq C(K) |A|,
\]
where $d(K)$ and $C(K)$ depend only on $K$ (and not on $|A|$).

More precisely, Ruzsa's version \cite{Ruzsa1994} gives
$d \leq 2K^4 \log(2K)$
and $C(K) = \exp(K^4 \log K)$.
\end{theorem}

The proof strategy (following Ruzsa \cite{Ruzsa1994}) uses two main ingredients:
\begin{enumerate}[label=(\roman*)]
  \item A \emph{Freiman homomorphism} argument to lift $A$ to a higher-dimensional
        torus $\Z_N^d$ where structure is easier to detect.
  \item A \emph{geometry of numbers} argument (Minkowski's theorem) to find a
        large GAP containing the image of $A$.
\end{enumerate}

\begin{definition}[Freiman Homomorphism]
A map $\phi: A \to G'$ between subsets of abelian groups is a
\emph{Freiman $s$-homomorphism} if
\[
  a_1 + \cdots + a_s = b_1 + \cdots + b_s \implies
  \phi(a_1) + \cdots + \phi(a_s) = \phi(b_1) + \cdots + \phi(b_s)
\]
for all $a_i, b_i \in A$. It is a \emph{Freiman $s$-isomorphism} if the converse
also holds.
\end{definition}

\begin{proposition}\label{prop:freiman_iso}
If $A \subset \Z$ has $|A + A| \leq K|A|$ and $n \geq |A|^2$, then there exists
a Freiman $2$-isomorphism $\phi: A \to \Z_n$ with $|\phi(A) + \phi(A)| \leq K|\phi(A)|$.
\end{proposition}

The key geometric lemma, due to Freiman, identifies GAPs via lattice points.

\begin{lemma}[Freiman's Lemma]\label{lem:freiman_lemma}
Let $\Lambda \subset \R^d$ be a lattice of determinant $\det(\Lambda)$ and let
$S \subset \Lambda$ be a finite set with $|S + S| \leq K|S|$. Then $S$ is
contained in a $d$-dimensional proper GAP $P$ with
$|P| \leq (2K)^d |S|$.
\end{lemma}

\begin{proof}[Proof sketch]
Apply the Plünnecke--Ruzsa inequality to get $|dS - dS| \leq K^{2d}|S|$.
The set $dS - dS$ spans a sublattice $\Lambda' \subseteq \Z^d$ of index at most
$K^{2d}|S|^d / |S| = K^{2d}|S|^{d-1}$. Minkowski's second theorem then provides
$d$ linearly independent short vectors in $\Lambda'$, which form the generators of
the desired GAP.
\end{proof}

%% ============================================================
\section{Green--Ruzsa Generalization}\label{sec:green_ruzsa}
%% ============================================================

Green and Ruzsa \cite{GreenRuzsa2007} extended Freiman's theorem to general abelian
groups, where torsion elements require a more flexible notion of "progression."

\begin{definition}[Coset Progression]
A \emph{coset $d$-progression} in an abelian group $G$ is a set of the form
$H + P$ where $H \leq G$ is a finite subgroup and $P$ is a $d$-dimensional GAP in
$G/H$ lifted to $G$.
\end{definition}

\begin{theorem}[Green--Ruzsa Theorem \cite{GreenRuzsa2007}]\label{thm:green_ruzsa}
Let $G$ be an abelian group and $A \subset G$ a finite set with $|A + A| \leq K|A|$.
Then $A$ is contained in a coset $d$-progression $H + P$ with
\[
  d \leq d(K), \qquad |H + P| \leq C(K)|A|,
\]
where $d(K)$ and $C(K)$ depend only on $K$.
\end{theorem}

\begin{proof}[Proof sketch]
The torsion-free case reduces to Freiman's theorem. For groups with torsion, one
first identifies the "torsion subgroup" $H$ generated by elements of small order
that appear in $A + A$. The quotient $G/H$ is torsion-free (in the relevant range),
and the image of $A$ in $G/H$ satisfies a doubling condition. Freiman's theorem
applied in $G/H$ gives a GAP, which lifts to a coset progression in $G$.
\end{proof}

\begin{remark}
For $G = \F_2^n$ (which is entirely $2$-torsion), the result gives a coset of a
subspace of bounded codimension, which is a much stronger structural statement
exploited in the Bogolyubov--Chang--Sanders line of work.
\end{remark}

%% ============================================================
\section{Bogolyubov's Method}\label{sec:bogolyubov}
%% ============================================================

Bogolyubov's lemma is a Fourier-analytic tool that converts the sumset assumption
into density information on a structured set.

\begin{lemma}[Bogolyubov's Lemma \cite{Bogolyubov1939}]\label{lem:bogolyubov}
Let $A \subseteq \Z_N$ with $|A| = \alpha N$. Then $2A - 2A$ contains a Bohr set
\[
  \mathrm{Bohr}(S, \rho) := \{n \in \Z_N \mid \|\xi n / N\| \leq \rho \text{ for all } \xi \in S\}
\]
with $|S| = O(\alpha^{-2})$ and $\rho = 1/4$.
\end{lemma}

\begin{proof}
Let $f = \mathbf{1}_A - \alpha$ (the balanced function of $A$). Its Fourier
transform satisfies $\hat{f}(0) = 0$ and $\|\hat{f}\|_2^2 = \alpha(1-\alpha)N$.
The convolution $f * f * f * f$ (which detects $2A - 2A$) is bounded below using
the large Fourier coefficients. Setting $S = \{\xi : |\hat{f}(\xi)| \geq \alpha^2 N / 2\}$
gives $|S| = O(\alpha^{-2})$ by Parseval, and the Bohr set argument completes
the proof.
\end{proof}

\begin{corollary}
If $|A + A| \leq K|A|$ in $\Z$, then $2A - 2A$ contains a long arithmetic
progression of length $\exp(-O(K^2)) |A|$.
\end{corollary}

%% ============================================================
\section{Connection to Roth's Theorem}\label{sec:roth}
%% ============================================================

Freiman's theorem and Bogolyubov's lemma have a deep connection to Roth's theorem
on 3-term arithmetic progressions.

\begin{theorem}[Roth's Theorem \cite{Roth1953}]\label{thm:roth}
Any subset $A \subseteq \{1, \ldots, N\}$ with
$|A| \geq c N / \log\log N$
contains a 3-term arithmetic progression.
\end{theorem}

The connection is as follows: if $A$ has no 3-term progressions, then $A$ has
small additive energy $E(A) := |\{(a_1, a_2, a_3, a_4) \in A^4: a_1+a_2=a_3+a_4\}|$,
which (via the Cauchy--Schwarz inequality) forces $|A + A|$ to be large relative
to $|A|$. Conversely, a Freiman-structured set (lying in a GAP) has many 3-term
progressions, which is the key tension exploited in density increment arguments.

\begin{proposition}[Additive Energy and Doubling]\label{prop:energy_doubling}
For $A \subseteq \Z$:
\[
  \frac{|A|^4}{|A + A|} \leq E(A) \leq |A|^2 \min(|A|, |A+A|)^{1/2}.
\]
The lower bound follows from Cauchy--Schwarz applied to the representation function
$r_{A+A}(n) = |\{(a, b): a+b=n\}|$.
\end{proposition}

%% ============================================================
\section{Sum-Free Sets}\label{sec:sumfree}
%% ============================================================

\begin{definition}[Sum-Free Set]
A set $A \subseteq G$ is \emph{sum-free} if $(A + A) \cap A = \emptyset$, i.e.,
there are no solutions to $a + b = c$ with $a, b, c \in A$.
\end{definition}

\begin{theorem}[Schur's Theorem \cite{Schur1916}]\label{thm:schur}
For every $k \geq 1$ there exists $N = S(k)$ (the $k$-th Schur number) such that
any $k$-coloring of $\{1, \ldots, N\}$ contains a monochromatic solution to
$a + b = c$.
\end{theorem}

\begin{proposition}[Density of Sum-Free Sets]\label{prop:sumfree_density}
Any sum-free subset $A \subseteq \{1, \ldots, N\}$ satisfies $|A| \leq N/3 + O(1)$.
The extremal example is the set of integers in $(N/3, 2N/3]$.
\end{proposition}

\begin{proof}
Consider the triples $(a, b, a+b)$ with $a, b, a+b \in \{1, \ldots, N\}$.
Each element of $A$ can play the role of $a+b$ in at most $|A \cap \{1,\ldots,x/2\}|$
triples, and similarly for $a$ and $b$. A counting argument shows $|A| \leq N/3$.
\end{proof}

The Plünnecke--Ruzsa inequality implies that sum-free sets cannot have too small
a doubling constant:

\begin{corollary}
A sum-free set $A \subseteq \Z$ satisfies $|A + A| \geq 2|A|$, hence
$\sigma[A] \geq 2$, with equality iff $A$ is an arithmetic progression of odd
common difference contained in an interval $(N/3, 2N/3]$.
\end{corollary}

%% ============================================================
\section{Sanders' Quasi-Polynomial Improvement}\label{sec:sanders}
%% ============================================================

The quantitative aspects of Freiman's theorem have been significantly sharpened.
The key question is: what is the optimal form of $C(K)$ in Theorem~\ref{thm:freiman}?

\begin{theorem}[Bilu's Theorem \cite{Bilu1999}]\label{thm:bilu}
In Freiman's theorem, one can take $|P| \leq \exp(O(K^2)) |A|$ with $d \leq K^2$.
\end{theorem}

\begin{theorem}[Sanders' Quasi-Polynomial Bound \cite{Sanders2012}]\label{thm:sanders}
If $A \subseteq \Z$ with $|A + A| \leq K|A|$, then $A$ is contained in a
$d$-dimensional GAP $P$ with
\[
  d \leq O(\log K), \qquad |P| \leq \exp(O(\log^4 K)) |A|.
\]
In particular, $C(K) = \exp(O(\log^4 K))$ is quasi-polynomial in $K$.
\end{theorem}

Sanders' proof uses a sophisticated iteration of Bogolyubov's lemma combined with
additive combinatorics arguments in Bohr sets. The key technical improvement is
handling the "density increment" in a Bohr set more efficiently.

\begin{conjecture}[Polynomial Freiman--Ruzsa Conjecture]\label{conj:pfr}
In Freiman's theorem, one can take $C(K) = K^{O(1)}$, i.e., the size of the
containing GAP is polynomial in $K$ times $|A|$.
\end{conjecture}

\begin{remark}
In March 2023, Gowers, Green, Manners and Tao \cite{GowersTao2023} announced a
proof of the Polynomial Freiman--Ruzsa (PFR) conjecture over $\F_2^n$, confirming
that over the boolean group any set with $|A+A| \leq K|A|$ is covered by $K^{O(1)}$
cosets of a subspace. Whether this implies the integer version remains open.
\end{remark}

%% ============================================================
\section{Inverse Sumset Theory}\label{sec:inverse}
%% ============================================================

Inverse sumset theory asks: given information about $|A + B|$, what can be said
about the structure of $A$ and $B$?

\begin{theorem}[Freiman--Ruzsa Theorem for Pairs]\label{thm:freiman_ruzsa_pairs}
Let $A, B \subseteq \Z$ with $|A + B| \leq K \sqrt{|A||B|}$. Then there exists a
$d$-dimensional GAP $P$ with $A \cup B \subseteq P$ and $|P| \leq C(K)|A+B|$,
for $d$ and $C(K)$ depending only on $K$.
\end{theorem}

\begin{conjecture}[Inverse Sumset Conjecture for General Groups]\label{conj:inverse_general}
For every abelian group $G$ and finite $A \subseteq G$ with $|A + A| \leq K|A|$,
the set $A$ is covered by $K^{O(1)}$ translates of a subgroup (or coset progression
of dimension $O(\log K)$) of size at most $K^{O(1)} |A|$.
\end{conjecture}

%% ============================================================
\section{The Polynomial Freiman--Ruzsa Conjecture over $\F_2^n$}\label{sec:pfr_f2}
%% ============================================================

Over the field $\F_2^n$, the analogue of a GAP is a coset of a subspace.

\begin{definition}[Covering Number]
The \emph{$K$-covering number} $\mathrm{cov}_K(A)$ of $A \subseteq \F_2^n$ is the
minimum size of a subspace $V$ such that $A$ is covered by at most $K$ cosets of $V$.
\end{definition}

\begin{theorem}[PFR over $\F_2^n$ \cite{GowersTao2023}]\label{thm:pfr_f2}
If $A \subseteq \F_2^n$ with $|A + A| \leq K|A|$, then $A$ is covered by
$2K^{12}$ cosets of a subspace $V$ with $|V| \leq 2K^{14}|A|$.
\end{theorem}

\begin{proof}[Proof outline]
The proof proceeds by a careful entropy-increment argument. Define the
\emph{Ruzsa distance} $d(A, B) = \log(|A-B| / \sqrt{|A||B|})$. The key step is
showing that if $d(A, A)$ is small, then $A$ is close (in Ruzsa distance) to a
subspace. The entropy method, combined with the \emph{fibring trick} and a
dimension-reduction argument, yields the polynomial bound.
\end{proof}

\begin{conjecture}[Optimal PFR Constant]\label{conj:pfr_optimal}
The optimal constant in Theorem~\ref{thm:pfr_f2} is $K^{O(1)}$ with the
exponent being at most $O(1)$ (specifically, conjectured to be $\leq 6$).
\end{conjecture}

%% ============================================================
\section{Quantitative Bounds: Summary Table}\label{sec:summary}
%% ============================================================

\begin{center}
\begin{tabular}{lll}
\toprule
\textbf{Result} & \textbf{Bound on $|P|/|A|$} & \textbf{Dimension $d$} \\
\midrule
Freiman (1966) & $\exp(\exp(K))$ & $\exp(K)$ \\
Ruzsa (1994)   & $\exp(K^4 \log K)$ & $2K^4 \log(2K)$ \\
Bilu (1999)    & $\exp(O(K^2))$ & $O(K^2)$ \\
Chang (2002)   & $\exp(O(K^2))$ & $O(K)$ \\
Sanders (2012) & $\exp(O(\log^4 K))$ & $O(\log K)$ \\
PFR over $\F_2^n$ (2023) & $K^{14}$ (polynomial) & --- \\
\bottomrule
\end{tabular}
\end{center}

%% ============================================================
\section{Open Problems}\label{sec:open}
%% ============================================================

We list the main open problems in the area.

\begin{conjecture}[Polynomial Freiman--Ruzsa over $\Z$]\label{conj:pfr_z}
There exist absolute constants $c, C > 0$ such that if $A \subset \Z$ with
$|A+A| \leq K|A|$, then $A$ is contained in a GAP $P$ with $|P| \leq K^C |A|$
and dimension $d \leq c \log K$.
\end{conjecture}

\begin{conjecture}[Optimal Doubling Constant]\label{conj:optimal_doubling}
For arithmetic progressions $P$ of length $n$, the doubling constant satisfies
$\sigma[P] = 2 - 1/|P|$. It is conjectured that the sets achieving $\sigma[A]
= 2 - 1/|A| + \varepsilon$ for small $\varepsilon$ are always "close" to arithmetic
progressions, but a precise stability theorem with optimal constants is open.
\end{conjecture}

\begin{conjecture}[Freiman--Bilu Exponent]\label{conj:freiman_bilu}
The correct exponent $\alpha$ in $|P| \leq K^\alpha |A|$ in
Conjecture~\ref{conj:pfr_z} satisfies $1 \leq \alpha \leq 2$. The lower bound
$\alpha \geq 1$ is trivially necessary; the upper bound $\alpha \leq 2$ is
conjectured based on random set constructions.
\end{conjecture}

\begin{conjecture}[Inverse Sumset for Non-Abelian Groups]\label{conj:nonabelian}
An appropriate analogue of Freiman's theorem holds for non-abelian groups
(with the role of GAPs played by cosets of approximate groups in the sense
of Hrushovski and Breuillard--Green--Tao).
\end{conjecture}

\begin{remark}
Breuillard, Green, and Tao \cite{BreuillardGreenTao2012} have proved a version of
Freiman's theorem for linear groups over fields, showing that sets of small doubling
in $GL_n(k)$ are essentially "nilpotent coset progressions."
\end{remark}

%% ============================================================
\section{Conclusion}
%% ============================================================

Freiman's theorem stands at the interface of additive combinatorics, harmonic
analysis, and ergodic theory. Starting from the qualitative insight that small
doubling implies arithmetic structure, the quantitative theory has advanced
dramatically: from Ruzsa's exponential bound (1994) to Sanders' quasi-polynomial
improvement (2012) to Gowers--Green--Manners--Tao's polynomial result over $\F_2^n$
(2023). The central open problem---whether a polynomial bound holds over $\Z$---remains
one of the most challenging and important questions in the field.

The Plünnecke--Ruzsa inequality, Bogolyubov's method, and the connection to Roth's
theorem on arithmetic progressions form a toolkit that has found applications far
beyond its original context, including in analytic number theory, theoretical computer
science, and the study of expanders.

%% Literaturverzeichnis
\begin{thebibliography}{99}

\bibitem{Bogolyubov1939}
N.\,N.~Bogolyubov,
\textit{On some arithmetical properties of almost periodic functions},
Ann. Chaire Math. Phys. Kiev \textbf{4} (1939), 185--205.

\bibitem{Bilu1999}
Y.~Bilu,
\textit{Structure of sets with small sumset},
Ast\'erisque \textbf{258} (1999), 77--108.

\bibitem{BreuillardGreenTao2012}
E.~Breuillard, B.~Green, T.~Tao,
\textit{The structure of approximate groups},
Publ.\ Math.\ IHES \textbf{116} (2012), 115--221.

\bibitem{Freiman1966}
G.\,A.~Freiman,
\textit{Foundations of a Structural Theory of Set Addition},
Kazan' Gos. Ped. Inst., Kazan', 1966 (in Russian);
English transl.\ AMS, Providence, 1973.

\bibitem{GreenRuzsa2007}
B.~Green, I.~Z.~Ruzsa,
\textit{Freiman's theorem in an arbitrary abelian group},
J.\ Lond.\ Math.\ Soc.\ \textbf{75} (2007), 163--175.

\bibitem{GowersTao2023}
W.\,T.~Gowers, B.~Green, F.~Manners, T.~Tao,
\textit{On a conjecture of Marton},
preprint arXiv:2311.05762 (2023).

\bibitem{Plunnecke1970}
H.~Pl\"unnecke,
\textit{Eine zahlentheoretische Anwendung der Graphentheorie},
J.\ reine angew.\ Math.\ \textbf{243} (1970), 171--183.

\bibitem{Roth1953}
K.\,F.~Roth,
\textit{On certain sets of integers},
J.\ Lond.\ Math.\ Soc.\ \textbf{28} (1953), 104--109.

\bibitem{Ruzsa1989}
I.\,Z.~Ruzsa,
\textit{An application of graph theory to additive number theory},
Scientia (Budapest) \textbf{3} (1989), 97--109.

\bibitem{Ruzsa1994}
I.\,Z.~Ruzsa,
\textit{Generalized arithmetical progressions and sumsets},
Acta Math.\ Hungar.\ \textbf{65} (1994), 379--388.

\bibitem{Ruzsa2009}
I.\,Z.~Ruzsa,
\textit{Sumsets and structure},
in: Combinatorial Number Theory and Additive Group Theory,
Birkh\"auser, 2009, 87--210.

\bibitem{Sanders2012}
T.~Sanders,
\textit{On the Bogolyubov--Ruzsa lemma},
Analysis \& PDE \textbf{5} (2012), 627--655.

\bibitem{Schur1916}
I.~Schur,
\textit{\"Uber die Kongruenz $x^m + y^m \equiv z^m \pmod{p}$},
Jahresber.\ Dtsch.\ Math.-Ver.\ \textbf{25} (1916), 114--117.

\bibitem{TaoVu2006}
T.~Tao, V.~Vu,
\textit{Additive Combinatorics},
Cambridge University Press, Cambridge, 2006.

\end{thebibliography}

\end{document}
